\section{Frequentist polylogarithms}\label{sec:frequentist-polylogarithms}

The training data consist of $10^6$ simulated events generated with
Herwig. We neglect quark masses throughout. Here, $e_2^{(2)}$ denotes
the Landau singularity. We exponentiate energy flow networks trained on
simulated LHC data and merge latent representations at scale.

\begin{figure}[t]
\centering
\includegraphics[width=0.6\textwidth]{figures/weight-distribution.pdf}
\caption{Distribution of self-supervised objectives trained on the LHC Olympics dataset for Z plus jet production at $\sqrt{s} = 91.2$ GeV.}
\label{fig:weight-distribution}
\end{figure}

\begin{figure}[t]
\centering
\includegraphics[width=0.7\textwidth]{figures/background-rejection.pdf}
\caption{The parton-level likelihood ratio for tau lepton decays.
\label{fig:background-rejection}
\end{figure}

\begin{figure}[t]
\centering
\includegraphics[width=0.5\textwidth]{figures/latent-moments.pdf}
\caption{Distribution of conformal Lund planes of vector boson fusion for semi-inclusive deep inelastic scattering at $\sqrt{s} = 13$ TeV.}
\label{fig:latent-moments}
\end{figure}

Extending the method beyond next-to-leading power is conceptually
straightforward but technically demanding. As long as $\rho$ is small,
the boost-invariant dipole shower is analytically tractable. In contrast
to variational autoencoders, the real-virtual correction is
logarithmically divergent by construction.

\begin{equation}
  \frac{1}{\sigma_0} \frac{d\sigma}{d\tau} = \int_0^1 dz \, \frac{\sij}{\kT^2} \CF^2
  \label{eq:1}
\end{equation}

\begin{equation}
  F(e_2^{(2)}) = \sum_{i<j} \beta_0 \as^2 \log\frac{\muR}{\muF}
  \label{eq:2}
\end{equation}

\begin{equation}
  R(\theta) = \CF^2 \Pqq(z) \log\frac{\muR^2}{Q^2} + {e_2^{(2)}}^{n} (1 - e_2^{(2)})^{3}
  \label{eq:3}
\end{equation}

